% No preamble. Notebook 16 classical coarse-action module. Unique labels: ca-.
% Place after the wr- response module. Standard theorem environments and \norm.

\section{The classical coarse action and its gauge identity}
\label{ca-section}

We remain on finite compatible periodic regulators, with fine
spacing $\eta=L^{-j}$ and final coarse spacing one. The action
$\mathcal S_\eta(U)$ is the coupling-free Wilson action in the
normalization of Notebook 15. Set
$\mathcal W_\eta(A)=\mathcal S_\eta(e^{i\eta A})$.
All coarse pairings below use the normalized color trace and
the physical coarse measure, which is one at this scale.
Transposes are bilinear after complexification. The constants
are independent of refinement and periodic volume.

Let $A_*(B)$ be the corrected-constraint minimizer in the fixed
flat gauge slice, and put $U_*(B)=e^{i\eta A_*(B)}$. Its
ordinary-feasible representative $U_{\rm ord}(B)$ is related
to $U_*(B)$ by the genuine gauge transformation constructed in
Notebook 15. Define
\begin{equation}\label{ca-definition}
 V(B)=\mathcal W_\eta(A_*(B))
     =\mathcal S_\eta(U_*(B))
     =\mathcal S_\eta(U_{\rm ord}(B)).
\end{equation}
These are equal holomorphic functions on a uniform complex
coarse supremum ball. Indeed the analytic Wilson density,
written with $P+P^{-1}$, is invariant under complex matrix
gauge conjugation as well. Only the real fields have the
minimum interpretation. The scalar total action scales with
volume; the uniform estimates below state that dependence.

\subsection{The multiplier is the exact derivative of the action}
\label{ca-envelope-section}

Write $\mathcal F(A)=QA+\mathcal R(A)$ and use the exact
chart and flat-slice space
\[
 A=\xi+H[B-\mathcal R(A)],\qquad
 \xi\in\mathcal X=\ker Q\cap\ker(R_0D^T).
\]
Here $R_0$ is the fixed scalar gauge projection, distinct from
the nonlinear remainder $\mathcal R$ and denoted $R$ in the
response module. Put
\begin{equation}\label{ca-chart-differential}
 T_A=(I+HD\mathcal R(A))^{-1},\qquad H_A=T_AH.
\end{equation}
Differentiating the analytic minimizer in a coarse direction $c$
gives the exact formula
\begin{equation}\label{ca-response-decomposition}
 D_BA_*(B)[c]=H_Ac+T_AD_B\xi(B)[c],\qquad A=A_*(B).
\end{equation}
The second summand is a slice tangent because
$D_B\xi(B)[c]\in\mathcal X$.

\begin{proposition}[Classical envelope identity]\label{ca-envelope}
Let $E(A)=D\mathcal W_\eta(A)$ denote the Euler field under
the physical fine pairing. The multiplier from the exact
stationarity construction satisfies
\begin{equation}\label{ca-first-derivative}
 DV(B)[c]=\langle m(B),c\rangle_1,\qquad
 m(B)=H_A^TE(A_*(B)).
\end{equation}
Equivalently,
\begin{equation}\label{ca-multiplier}
 \begin{split}
 m(B)={}&(I+H^TD\mathcal R(A_*)^T)^{-1}\\
 &\quad\cdot\{q[B-\mathcal R(A_*)]+H^TN(A_*)\},
 \end{split}
\end{equation}
where $q=q_M$ is the physical flat Maxwell precision and
$E(A)=KA+N(A)$ is the ordinary-coordinate force.
\end{proposition}

\begin{proof}
Pair \eqref{ca-response-decomposition} with $E(A_*)$.
Slice stationarity makes
$\langle E(A_*),T_AD_B\xi(B)[c]\rangle_\eta=0$.
The remaining term is exactly \eqref{ca-first-derivative}.
This argument needs only slice stationarity, not a separate
full-constraint multiplier theorem or gauge coverage.

The transpose push-through identity gives
$H_A^T=(I+H^TD\mathcal R(A_*)^T)^{-1}H^T$.
The flat identities $H^TK\xi=0$ and $H^TKH=q$ imply
$H^TKA_*=q[B-\mathcal R(A_*)]$, proving
\eqref{ca-multiplier}. The force term is evaluated as
$H^TN(A_*)=(S_{\rm jet}H)^Tf_\eta(S_{\rm jet}A_*)$;
no rough force coefficient is differentiated. All identities
extend bilinearly to the complex ball.
\end{proof}

Put $b=\norm{B}_\infty$. The uniform estimates
$A_*=HB+O(b^2)$, $\mathcal R(A_*)=O(b^2)$,
$H^TN(A_*)=O(b^2)$ and
$\norm{D\mathcal R(A_*)^T}_{\infty\to\infty}=O(b)$
in \eqref{ca-multiplier} prove
\begin{equation}\label{ca-m-near-q}
 \norm{m(B)-qB}_\infty\le Cb^2,
 \qquad m(0)=0,\quad Dm(0)=q.
\end{equation}
The multiplier is holomorphic on a uniform coarse ball.
In particular $D^2V(0)=q$. Its unweighted estimate
$\norm{Dm(B)-q}_{\infty\to\infty}\le Cb$ also follows
directly by Cauchy: for $b>0$ use a direction circle of radius
$b$ inside a smaller common ball, and apply
\eqref{ca-m-near-q} at points of norm at most $2b$.
At $b=0$ use the displayed Taylor coefficient.

Let $N_b$ denote the number of independent positive coarse
bonds. Since $V(0)=0$, radial integration gives
\begin{equation}\label{ca-radial}
 V(B)=\frac12\langle B,qB\rangle_1
       +\int_0^1\langle m(tB)-tqB,B\rangle_1\,dt.
\end{equation}
The fixed normalized color trace bound and
\eqref{ca-m-near-q} therefore yield
\begin{equation}\label{ca-cubic-remainder}
 \boxed{\quad
 \left|V(B)-\frac12\langle B,qB\rangle_1\right|
          \le C N_b\norm{B}_\infty^3.
 \quad}
\end{equation}
This holds on the complex ball. It is a volume-uniform
density estimate, not a volume-independent total-action bound.

\subsection{Local coarse gauge invariance}
\label{ca-gauge-section}

For a real coarse site field $\omega$ put $g=e^{i\omega}$,
$D=e^{iB}$ and
\begin{equation}\label{ca-coarse-gauge}
 D^g(c)=g(c_-)D(c)g(c_+)^{-1},\qquad
 B^g=\frac1i\log D^g.
\end{equation}
All logarithms use the common small branch. At coarse spacing
one the local analytic multiplication gives a constant
$C_c\ge1$ such that
\begin{equation}\label{ca-coarse-gauge-bound}
 \norm{B^g}_\infty\le C_c(b+w),\qquad
 w=\norm{\omega}_\infty.
\end{equation}
No additional regularity is required of $\omega$: coarse
first differences are bounded by $2w$.

Let $b_0,r_0$ be the ordinary-minimum radii of Notebook 15,
and write
$\norm{A_{\rm ord}(B)}_{N_\alpha}\le C_Ab$.
The smooth point-value extension there satisfies
\begin{equation}\label{ca-smooth-gauge-lift}
 (E_{\rm pt}\omega)|_{\rm coarse}=\omega,\qquad
 \norm{E_{\rm pt}\omega}_{C^{2,\alpha}}\le C_Ew.
\end{equation}
Set $h=e^{iE_{\rm pt}\omega}$, so $h|_{\rm coarse}=g$.
The fine gauge-action bound, valid for both $h$ and $h^{-1}$,
is
\begin{equation}\label{ca-fine-gauge-bound}
 \norm{A^h}_{N_\alpha}
       \le C_G(\norm A_{N_\alpha}+C_Ew).
\end{equation}

Choose one $\delta>0$ so that all logarithms and gauge actions
stay in their established domains and
\begin{equation}\label{ca-comparison-margins}
 C_c\delta<b_0/2,\qquad
 C_G(C_AC_c+C_E)\delta<r_0/2.
\end{equation}
For real $b+w<\delta$, ordinary gauge covariance makes
$U_{\rm ord}(B)^h$ feasible for coarse datum $B^g$.
Its potential norm is below $r_0$, so minimality at $B^g$
and exact gauge invariance give
\[
 V(B^g)\le\mathcal S_\eta(U_{\rm ord}(B)^h)=V(B).
\]
For the reverse comparison use
$U_{\rm ord}(B^g)^{h^{-1}}$. Its ordinary coarse average is
$(D^g)^{g^{-1}}=D$, and its norm is bounded by
\[
 C_G\{C_AC_c(b+w)+C_Ew\}<r_0.
\]
It is therefore an admissible competitor at $B$, proving
$V(B)\le V(B^g)$. We have proved
\begin{equation}\label{ca-gauge-invariance}
 \boxed{\qquad V(B^g)=V(B)\qquad(b+w<\delta).\qquad}
\end{equation}
Both comparisons lie strictly inside the same uniform minimum
radius. The reverse one uses the separately controlled minimizer
at $B^g$, not an assumed inverse of the first comparison map.

The two sides of \eqref{ca-gauge-invariance} are holomorphic
in $(B,\omega)$ on a uniform smaller complex product ball.
Their equality on an open set of the full real form implies
equality on that connected product ball by the finite-dimensional
holomorphic identity theorem, applied coordinate by coordinate.
For example, reduce $\delta$ to respect all analytic domains
and take $\norm B_\infty,\norm\omega_\infty<\delta/3$.
This extension is an analytic identity; it is not a minimum
claim for complex fields.

\subsection{The exact classical Ward identity}
\label{ca-ward-section}

Define the coarse left exponential differential and the
coarse covariant scalar difference by
\begin{equation}\label{ca-left-differential}
 \mathcal J_Bv=\int_0^1\operatorname{Ad}(e^{itB})v\,dt,
 \qquad
 (D_D\omega)(c)=\operatorname{Ad}_{D(c)}\omega(c_+)-\omega(c_-).
\end{equation}
The operator $\mathcal J_B$ is invertible on the fixed small
complex ball. The ordinary coarse-potential velocity under
$g(t)=e^{it\omega}$ is
\begin{equation}\label{ca-gauge-velocity}
 \left.\frac{d}{dt}B^{g(t)}\right|_{t=0}
             =-\mathcal J_B^{-1}D_D\omega.
\end{equation}
The coarse spacing is one, so there is no additional length
factor. Differentiate \eqref{ca-gauge-invariance} and use
\eqref{ca-first-derivative}:
\[
 0=-\langle m(B),\mathcal J_B^{-1}D_D\omega\rangle_1
   =-\langle D_D^T\mathcal J_B^{-T}m(B),\omega\rangle_1.
\]
The site field $\omega$ is arbitrary. Hence
\begin{equation}\label{ca-ward}
 \boxed{\qquad D_D^T\mathcal J_B^{-T}m(B)=0.\qquad}
\end{equation}
The real identity extends holomorphically. Thus the
left-coordinate classical coarse current is
$j_{\rm cl}(B)=\mathcal J_B^{-T}m(B)$, not generally $m(B)$.
No quantum determinant or quantum Ward assertion is involved.

\subsection{A local classical Hessian and an auxiliary positive inverse}
\label{ca-precision-section}

The weighted multiplier response \eqref{wr-dm-near-q}, together
with the envelope identity, gives
\begin{equation}\label{ca-hessian-weighted}
 Q_{\rm cl}(B):=D^2V(B)=Dm(B),\qquad
 \norm{[Q_{\rm cl}(B)-q]c}_{\infty,w}
       \le C_\kappa b\norm c_{\infty,w}
\end{equation}
for every admissible weight in \eqref{wr-weight}. As the
Hessian of a holomorphic scalar function, $Q_{\rm cl}(B)$
is symmetric for the physical bilinear transpose.
It differentiates in the ordinary coarse potential $B$.
A change to a based left-exponential fluctuation chart
requires the coordinate Jacobians and its background-current
second-derivative term; that differently based Hessian is
not being identified with $Q_{\rm cl}(B)$ here.

Use a fixed trace-orthonormal color and orientation basis in
each coarse cell. For a block kernel $T$ define the weighted
Schur norm
\begin{equation}\label{ca-schur}
 \norm T_\kappa=\max\left\{
   \sup_n\sum_m e^{\kappa d(n,m)}\norm{T_{nm}},\,
   \sup_m\sum_n e^{\kappa d(n,m)}\norm{T_{nm}}\right\}.
\end{equation}
Take $w_n(m)=e^{-\kappa d(n,m)}$ in
\eqref{ca-hessian-weighted}. In each scalar output row,
choose input phases to match its matrix entries. This bounds
its absolute row sum with the factor $e^{\kappa d(n,m)}$.
Summing the fixed number of component rows converts this to
matrix-block norms; transpose symmetry gives the column bound.
Consequently
\begin{equation}\label{ca-hessian-schur}
 \norm{Q_{\rm cl}(B)-q}_\kappa\le C'_{\kappa}b.
\end{equation}
The component factor depends only on dimension and color,
not on refinement or volume. These finite-regulator row tests
are bounded directions even for decaying weights.

Let $a$ lie in a fixed interval $0<a_-\le a\le a_+$.
The flat baseline estimates give, after fixing a sufficiently
small common $\kappa$,
\begin{equation}\label{ca-baseline-inverse}
 \norm{(q+aI)^{-1}}_\kappa\le r_\kappa
\end{equation}
uniformly. Set $P_{\rm cl}(B)=Q_{\rm cl}(B)+aI$.
For
\begin{equation}\label{ca-inverse-radius}
 C'_\kappa b\le a_-/2,\qquad
 r_\kappa C'_\kappa b\le1/2,
\end{equation}
the real positive baseline $q\ge0$ and the Schur operator
bound imply
\begin{equation}\label{ca-accretivity}
 \operatorname{Re}\langle x,P_{\rm cl}(B)x\rangle_{\rm Herm}
                 \ge\frac{a_-}{2}\norm x_2^2
\end{equation}
even at complex $B$. The numerical-range pairing here is
Hermitian; it is distinct from the analytic bilinear transpose.
The Schur norm is a matrix algebra norm, so the Neumann series
also gives
\begin{equation}\label{ca-inverse-bounds}
 \begin{split}
 \norm{P_{\rm cl}(B)^{-1}}_\kappa&\le2r_\kappa,\\
 \norm{P_{\rm cl}(B)^{-1}-(q+aI)^{-1}}_\kappa
                         &\le2r_\kappa^2C'_\kappa b.
 \end{split}
\end{equation}
This inverse is holomorphic and exponentially local in
coarse-cell distance, with uniform constants.

The positive $a$ is an auxiliary penalty added to a coarse
classical Hessian. This is not a physical Yang--Mills mass gap
or the precision of an integrated quantum effective action.
Spatial kernel decay does not assert strict dependence on
bounded background supports. Fluctuation determinants,
Jacobian factors, quantum interactions and their iteration
remain separate tasks. No finite total action is assigned
here to arbitrary bounded fields on the infinite lattice.
